\section{Passage path: can it move through both chambers and then work?}
\label{sec:passage}

The eighth question closes the loop. Even a good bill needs a path, and a member
wants to know it can both pass and then be implemented without embarrassment.

The path is ordinary regular order, run in parallel across the chambers. The bill is
introduced and referred to committee, refined in subcommittee and full-committee
markup, and passed on the House floor; a Senate companion follows the same sequence;
and the two versions are reconciled by amendment exchange or conference before the
bill is enrolled and presented \cite{congress2024lawsmade}. The thresholds are the
familiar ones: a simple majority of the House, sixty votes to invoke cloture in the
Senate, and a simple majority to pass it. The appendix figures trace this journey,
the two-chamber exchange, the markup as a branch-and-merge of the text, and the vote
thresholds as formal requirements.

Implementation is already scoped, which is what lets a member promise the bill will
work. Section 515D directs the Secretary to issue final regulations within 365 days,
takes effect on the first day of the first calendar quarter at least one year after
enactment, and gives existing enrolling investigations 180 days to comply, so the
runway is defined rather than open-ended \cite{Kawchak2026HR9510Bill}. More
important, the system the bill governs has already been demonstrated. In validated
simulation the National Platform treated 168 patients with 29 robots across 15 cancer
types at 99.7 percent uptime with zero patient harm events, and it is built to operate
as a fully automated sponsor across many sites \cite{Kawchak2026NationalPlatformPhysicalAI,
Kawchak2026FullyAutomatedSponsorPhysical}. \autoref{fig:capstone} closes the
framework by tracing one member's yes up to the statute and the platform ready to
implement it.

\begin{psgfig}
\node[psgone] (e) at (0,0) {Eight questions\\answered};
\node[psgevid] (y) at (4,0) {One member's\\yes};
\node[psgtwo] (r) at (8,0) {Markup and\\report};
\node[psgact] (p) at (12,0) {House and\\Senate passage};
\node[psgact] (l) at (5.2,-1.7) {Enacted\\section 515D};
\node[psgevid] (i) at (9.2,-1.7) {Platform ready\\to implement};
\draw[psgedge] (e)--(y); \draw[psgedge] (y)--(r); \draw[psgedge] (r)--(p);
\draw[psgedge] (p) to[out=-90,in=0,looseness=0.7] (i.east);
\draw[psgedge] (i)--(l);
\end{psgfig}
\vspace{-0.05cm}
\psgcaption{\textbf{Figure 9.} From one member's yes to enacted law and a platform
ready to implement it.}
\label{fig:capstone}

% \vspace{-0.2cm}

The passage path is answered with a plan: ordinary regular order in both chambers,
known vote thresholds, a defined implementation runway, and a platform already shown
to work, so a yes vote is a vote for something that will function.

\subsection*{Synthesis: each question and the provision that answers it}

The framework closes where it began, with the eight questions, now each paired with
the specific provision, score line, or evidence that settles it.
\autoref{tab:synthesis} is the one page a member can carry into a markup.

\needspace{15\baselineskip}\begin{table}[H]
\footnotesize
\begin{tabularx}{\textwidth}{@{}L{2.25cm} Y@{}}
\toprule
\textbf{Question} & \textbf{What answers it in H. R. 9510}\\
\midrule
Mandate & Law gaps (embodiment and sequencing, financial data, cost asymmetry), tied to preventable med. error.\\
Authority & Commerce-power amendment to FD\&C Act adding 515D, with construction rule and State savings clause.\\
Safety & Verification before generation: ten-gates bound to consensus standards, 3 hard catastrophe predicates.\\
Fiscal & 58 million dollar authorization, no net mandatory spending, no new fee, apx.\ 19x lower verification cost.\\
Constituents & A platform that scales to 50 through 100 sites across all regions, reaching rural and underserved districts.\\
Bipartisanship & One provision a deregulation-minded and a safety-minded member can each support on their own terms.\\
Coalition & Clinicians, patients, industry, and the States, with every standard objection routed to a cited answer.\\
Passage path & Ordinary regular order in both chambers, vote thresholds, implementation runway, demonstration.\\
\bottomrule
\end{tabularx}
\caption{The synthesis: each of the eight questions and the provision, score
line, or evidence in H. R. 9510.}
\label{tab:synthesis}
\end{table}

\clearpage